\documentclass{amsart}
\usepackage[a4paper,left=3cm,right=3cm,top=3cm,bottom=4cm]{geometry}
\usepackage{amsmath,amssymb,amsthm}
\usepackage{microtype}
\usepackage[colorlinks=true,linkcolor=blue,citecolor=blue,urlcolor=blue]{hyperref}

\newtheorem{theorem}{Theorem}
\DeclareMathOperator{\Aut}{Aut}
\DeclareMathOperator{\Sym}{Sym}
\DeclareMathOperator{\Hnm}{H}
\DeclareMathOperator{\Jac}{Jac}
\DeclareMathOperator{\rk}{rk}
\setlength{\parindent}{15pt}
\setlength{\parskip}{0pt plus 1pt}
\allowdisplaybreaks[1]

\title[The tangent Gauss map]{A counterexample to generic injectivity of the tangent Gauss map}
\author{Xiaoxiang Zhou}
\address{Institut f\"ur Mathematik, Humboldt-Universit\"at zu Berlin,
12489 Berlin, Germany}
\email{xiaoxiang.zhou@hu-berlin.de}
\date{November 3, 2026}

\begin{document}
\maketitle

We work over $\mathbb{C}$. For a curve $C\subset A$ in an abelian variety,
the tangent Gauss map sends a smooth point to its tangent line, translated
to the origin:
\[
\phi_C:C^{\mathrm{sm}}\longrightarrow\mathbb{P}(T_0A),\qquad
p\longmapsto d(t_{-p})_p(T_pC).
\]
We call $C$ nondegenerate if it is not contained in a translate of a
proper abelian subvariety.

\begin{theorem}\label{thm:weak}
There exists a nondegenerate curve in an abelian threefold whose tangent
Gauss map has generic degree $3$.
\end{theorem}

\section{Strategy}

Let $C$ be a smooth nonhyperelliptic curve of genus $4$, and let
$\mathcal{L}$ be a line bundle of degree $3$ with $h^0(\mathcal{L})=2$.
It is base-point-free, since a base point would give a degree-two pencil.
Thus $|\mathcal{L}|:C\longrightarrow\mathbb{P}^1$ has degree $3$.
We seek the following two properties:
\[
\mathcal{L}^{\otimes2}\cong\omega_C,\qquad
\Sym^2\Hnm^0(\mathcal{L})=q^*\Hnm^0(\Omega_A^1)
\quad\text{for a quotient }q:\Jac(C)\longrightarrow A.
\]
Here we identify $\Sym^2\Hnm^0(\mathcal{L})$ with its image in
$\Hnm^0(\omega_C)$ under multiplication. Its dimension is $3$: a relation
between the three products of a basis would force the ratio of its
sections to be constant. We call the second property the
\emph{cotangent-space condition}. The resulting linear system defines
\[
\phi:C\overset{|\mathcal{L}|}{\longrightarrow}\mathbb{P}^1
\overset{\nu_2}{\longrightarrow}\mathbb{P}^2,
\]
where $\nu_2$ is the quadratic Veronese embedding. We will realize
$\phi$ as a tangent Gauss map.

Start with a smooth hyperelliptic curve $C_0$ of genus $4$ admitting a
tangential cover $\pi_0:(C_0,p_0)\longrightarrow(E_0,0)$ at a Weierstrass
point. Such covers exist, for example among the Lam\'e spectral curves
\cite[Table~3 and Theorem~L]{maier} and \cite[Proposition~3.1]{eilbeck}.
Tangentiality means that $\pi_0^*E_0\subset\Jac(C_0)$ is tangent to the
Abel--Jacobi image at $p_0$. Put
$\mathcal{L}_0=\mathcal{O}_{C_0}(3p_0)$ and
$A_0=\Jac(C_0)/\pi_0^*E_0$, with quotient map $q_0$. Then
\[
\mathcal{L}_0^{\otimes2}\cong\omega_{C_0},\qquad
\Sym^2\Hnm^0(\mathcal{L}_0)
=\Hnm^0\left(\omega_{C_0}(-2p_0)\right)
=q_0^*\Hnm^0(\Omega_{A_0}^1).
\]
Indeed, the products of $1,x\in\Hnm^0(\mathcal{L}_0)$ span
$\Hnm^0\left(\mathcal{O}_{C_0}(4p_0)\right)$, where $x$ is a
hyperelliptic coordinate with a double pole at $p_0$. Also,
$\Hnm^0\left(\omega_{C_0}(-2p_0)\right)
=\Hnm^0\left(\omega_{C_0}(-p_0)\right)$ is the annihilator of the
tangent line at $p_0$. Conversely, if this annihilator comes from an
abelian quotient, its connected kernel is an elliptic curve tangent to
$C_0$; dualizing the inclusion and composing with the Abel--Jacobi map
gives a tangential cover.

\section{A dimension count for deformations}

Let $B$ be a small simply connected marked deformation space of the spin
curve $(C_0,\mathcal{L}_0)$. Its dimension is $3g-3=9$. The theta-null
locus
\[
S=\left\{t\in B\;\middle|\;h^0(\mathcal{L}_t)\geqslant2\right\}
\]
is locally defined by one theta constant. After shrinking $B$, upper
semicontinuity gives $h^0(\mathcal{L}_t)=2$ on $S$. Put
\[
W_t:=\Sym^2\Hnm^0(\mathcal{L}_t)\subset\Hnm^0(\omega_{C_t});
\]
these spaces form a holomorphic subbundle of rank $3$.

Set
\[
V_0=q_0^*\Hnm^1(A_0;\mathbb{Q})
\subset\Hnm^1(C_0;\mathbb{Q}),
\]
and let $V_t\subset\Hnm^1(C_t;\mathbb{Q})$ be its parallel transport
under the Gauss--Manin connection. To preserve the cotangent-space
condition, we use the following equivalent formulations:
\[
\begin{aligned}
&\begin{cases}
V_t\text{ remains a Hodge substructure},\\
V_{t,\mathbb{C}}\cap\Hnm^0(\omega_{C_t})=W_t
\end{cases}\\
\Longleftrightarrow\;& V_{t,\mathbb{C}}=W_t\oplus\overline{W}_t\\
\Longleftrightarrow\;& W_t\subseteq V_{t,\mathbb{C}}\\
\Longleftrightarrow\;&
\rk\left[W_t\longrightarrow\Hnm^1(C_t;\mathbb{C})/V_{t,\mathbb{C}}\right]=0.
\end{aligned}
\]
The middle equivalence follows from rationality of $V_t$ and Hodge type,
since $\dim W_t=3$ and $\dim V_t=6$. The last condition is the vanishing
of a bundle map of ranks $3$ and $2$, so it imposes at most six equations
on $S$. Consequently, every irreducible component $Z$ of its zero locus
through the seed satisfies
\[
\dim Z\geqslant9-1-6=2.
\]
On $Z$, the Hodge substructure $V_t$ gives an abelian quotient
$q_t:\Jac(C_t)\longrightarrow A_t$ with connected kernel and
$q_t^*\Hnm^0(\Omega_{A_t}^1)=W_t$.

We compare this with the hyperelliptic part. On a hyperelliptic curve
$C_t$ of genus $4$, a degree-three pencil must have a base point:
otherwise, together with the hyperelliptic pencil it would give a
birational map to a curve of bidegree $(2,3)$ in
$\mathbb{P}^1\times\mathbb{P}^1$, of arithmetic genus $2$.
Writing $\mathcal{L}_t=H_t(p_t)$, where $H_t$ is the hyperelliptic pencil,
the spin condition and $\omega_{C_t}\cong H_t^{\otimes3}$ give
$\mathcal{O}_{C_t}(2p_t)\cong H_t$. Hence $p_t$ is a Weierstrass point,
and
\[
\mathcal{L}_t\cong\mathcal{O}_{C_t}(3p_t),\qquad
W_t=\Sym^2\Hnm^0(\mathcal{L}_t)
=\Hnm^0\left(\omega_{C_t}(-2p_t)\right).
\]
By the tangency criterion above, these points give tangential covers
$\pi_t:(C_t,p_t)\longrightarrow(E_t,0)$. Their degree is locally fixed:
the integral homology lattice of the elliptic kernel of $q_t$ is the
annihilator of $V_t$, fixed under the marking, and dualizing its inclusion
fixes the cohomology map of $\pi_t$ and its cup-product pairing.

For a fixed elliptic curve and degree there are only finitely many
hyperelliptic tangential covers
\cite[Proposition~2.5, 3.; Section~2, Remark~1, 3.--4.]{treibich2025}.
Indeed, the relation $2n+1=\sum_{i=0}^3\mu_i^2+4d$ leaves finitely many
types for a fixed degree $n$, and each corresponding Severi set is
finite. Since elliptic curves have one modulus and the choices of
$p_t$ are finite,
\[
\dim\left(Z\cap\mathcal{H}_4\right)\leqslant1<2\leqslant\dim Z.
\]
We can therefore choose a nonhyperelliptic point of $Z$.

For this point, compose the Abel--Jacobi map based at $b\in C$ with $q$:
$f=q\circ j_b:C\longrightarrow A$. The pulled-back invariant forms
are the products of sections of $\mathcal{L}$, so $f$ is immersive and
its tangent map is $\phi$. Its image contains the origin and generates
$A$, hence is nondegenerate. If $f$ had degree $e>1$ onto the
normalization of its image, then $e\mid3$, since $\phi$ factors through
that normalization. Thus $e=3$, and the normalization would map with
degree $1$ to the conic $\nu_2(\mathbb{P}^1)$. This is impossible in an
abelian variety. Therefore $C$ is the normalization of $f(C)$ and the
tangent Gauss map of $f(C)$ has generic degree $3$, proving
Theorem~\ref{thm:weak}.

\section{Simplicity and symmetries}

\begin{theorem}\label{thm:strong}
The abelian threefold can be chosen simple. The resulting curve is
preserved by no nonidentity automorphism of $A$ with scalar differential.
\end{theorem}

First choose the hyperelliptic seed so that $A_0$ is simple; the existence
of such a seed is checked below. For every nonzero proper rational
subspace $R\subset V_0$ on which the polarization is nondegenerate, let
$e_R$ be the orthogonal projector. The condition
\[
e_R(W_t)\subset W_t
\]
defines a closed analytic subset of $Z$. It is proper because it misses
the seed, and it holds precisely when $R$ is a Hodge substructure.
Every abelian subvariety of $A_t$ gives such a substructure. There are
only countably many $R$, so Baire's theorem allows us to choose a point
outside these subsets and $\mathcal{H}_4$. Its quotient $A$ is simple.

Now suppose $\tau\in\Aut(A)$ preserves $f(C)$ and has scalar
differential, using translations to identify tangent spaces. It lifts
to an automorphism $\sigma$ of $C$ with $\phi\circ\sigma=\phi$.
If $\sigma$ is the identity, then $\tau$ fixes a generating curve
pointwise and is the identity. Otherwise, the order of $\sigma$ divides
$\deg\phi=3$, so it is $3$ and
$C/\langle\sigma\rangle\cong\mathbb{P}^1$. Riemann--Hurwitz gives a
fixed point of $\sigma$. Taking its image as the origin makes $\tau$
a group automorphism $u$ of order $3$. Its scalar differential is a
primitive cube root of unity $\lambda$.
Then $\mathbb{Q}(\lambda)$ acts on $\Hnm_1(A;\mathbb{Q})$, and the
complex structure on real homology is
$\left(u_*-\operatorname{Re}(\lambda)\right)/\operatorname{Im}(\lambda)$.
Every one-dimensional $\mathbb{Q}(\lambda)$-subspace therefore has a
complex-invariant realification and defines an elliptic subvariety of
$A$, contradicting simplicity. This proves the symmetry assertion.

\medskip
\noindent\emph{Existence of a seed with simple quotient.}
Take $E_0:y^2=4x^3-4x-1$, let $e_1,e_2,e_3$ be the roots of
$4e^3-4e-1$, and take the smooth projective model of
\[
C_0:\quad \nu^2=(B^3-208B+560)
\prod_{i=1}^3\left(B^2+10e_iB-35e_i^2-28\right).
\]
This is the genus-four Lam\'e spectral curve, with a degree-ten cover
of $E_0$ tangential at infinity
\cite[Table~3 and Theorem~L]{maier} and \cite[Proposition~3.1 and the discussion
after (17)]{eilbeck}.
The cover is defined over $\mathbb{Q}$, and $E_0$ has discriminant $37$.
Modulo $11$, the polynomial on the right is the square-free polynomial
$B^9+B^4+6B^3+9B^2+7B+5$; hence $C_0$ has genus $4$ and good reduction.
Direct counting gives
\[
\left(\#C_0(\mathbb{F}_{11^r})\right)_{r=1}^4
=(10,120,1312,15312),\qquad \#E_0(\mathbb{F}_{11})=7.
\]
Newton's identities and division by the elliptic factor $T^2-5T+11$
give the Frobenius polynomial of $A_0=\Jac(C_0)/\pi_0^*E_0$:
\[
h(T)=T^6+3T^5+5T^4-14T^3+55T^2+363T+1331.
\]
It is irreducible modulo $3$, and $11\nmid14$, so this reduction is
simple and ordinary. An ordinary simple abelian threefold over a
finite field of odd cardinality is absolutely simple unless its
Frobenius polynomial has the form $T^6+cT^3+q^3$
\cite[Proposition~3.4]{arias}. Thus this reduction is absolutely simple,
and so is $A_0$ in characteristic zero, since a proper abelian
subvariety would specialize to one after a finite extension.
% The existing verify_tangent_seed.py reproduces the reduction, point
% counts, Frobenius factorization, and irreducibility checks.

\begin{thebibliography}{9}
\raggedright
\bibitem{arias}
S.~Arias-de-Reyna, C.~Armana, V.~Karemaker, M.~Rebolledo,
L.~Thomas, and N.~Vila,
\emph{Large Galois images for Jacobian varieties of genus 3 curves},
\href{https://arxiv.org/abs/1507.05913}{arXiv:1507.05913}, 2016.

\bibitem{eilbeck}
J.~C.~Eilbeck, V.~Z.~Enolskii, and E.~Previato,
\emph{Varieties of elliptic solitons},
J.\ Phys.\ A: Math.\ Gen.\ \textbf{34} (2001), 2215--2227,
\href{https://doi.org/10.1088/0305-4470/34/11/314}{doi:10.1088/0305-4470/34/11/314}.

\bibitem{maier}
R.~S.~Maier,
\emph{Lam\'e polynomials, hyperelliptic reductions and Lam\'e band structure},
\href{https://arxiv.org/abs/math-ph/0309005v4}{arXiv:math-ph/0309005v4}, 2004.

\bibitem{treibich2025}
A.~Treibich,
\emph{Hyperelliptic tangential covers and even elliptic finite-gap
potentials, back and forth},
\href{https://arxiv.org/abs/2501.16483v1}{arXiv:2501.16483v1}, 2025.
\end{thebibliography}
\end{document}
